\clearpage
\VerdierPhysicalPageAnchor{0248}%
\begingroup\small\setlength{\parskip}{0pt}\setlength{\emergencystretch}{1.5em}
\setlength{\abovedisplayskip}{4pt plus 2pt minus 2pt}%
\setlength{\belowdisplayskip}{4pt plus 2pt minus 2pt}%

\VerdierSemanticAnchor{ch4d-4-6-8}
\noindent\textbf{命題 4.6.8.}\quad
\textit{$Y$ を $\mathcal A$ の複体とし、
$\varepsilon\colon Y\longrightarrow I^{\bullet\bullet}$ を $Y$ の
Cartan--Eilenberg 入射分解とする。$Y$ が
$\operatorname{comp}^{+}(\mathcal A)$ の対象であるか、または分解
$I^{\bullet\bullet}$ が有限長であるとき、単純複体
$\int_2 I^{\bullet\bullet}$ は入射型複体であり、複体の射}

\VerdierSemanticAnchor{ch4d-4-6-8-morphism}
\[
\textit{$\int_2\varepsilon\colon
Y\longrightarrow\int_2 I^{\bullet\bullet}$}
\]
\textit{は擬同型である。}

\VerdierSemanticAnchor{ch4d-4-6-8-proof-spectral-objects}
実際、射 $\varepsilon$ は、$Y$ の $D(\mathcal A)$ に値をとる第一スペクトル
対象 (4.1.5) と、第一次数でフィルター付けした $I^{\bullet\bullet}$ が定める
$D(\mathcal A)$ に値をとるスペクトル対象 ((4.1.7), (4.1.8)) との間の射を
誘導する。標準コホモロジー的関手
$H\colon D(\mathcal A)\longrightarrow\mathcal A$ は、これら二つのスペクトル
対象を、アーベル圏 $\mathcal A$ に値をとるスペクトル対象へ移す。
$\varepsilon$ がこれらのスペクトル対象の間の同型を誘導することを示せば
十分である。$Y$ が定めるスペクトル対象は常に定常であり（第 II 章、4.4.2）、
上の仮定のもとでは $I^{\bullet\bullet}$ が定めるスペクトル対象も定常である。
したがって、$\varepsilon$ が初期項上で同型を誘導することを示せば十分である。
これらの初期項は、一方では

\VerdierSemanticAnchor{ch4d-4-6-8-initial-y}
\[
\mathrm I_1^{p,q}=H^q(Y^p),
\]
であり、他方では

\VerdierSemanticAnchor{ch4d-4-6-8-initial-i}
\[
{}'\mathrm I_1^{p,q}=H^q(I^{p,\bullet}).
\]
である。さらに、$\varepsilon$ が初期項上で誘導する射は、複体の射

\VerdierSemanticAnchor{ch4d-4-6-8-epsilon-p}
\[
\varepsilon^p\colon Y^p\longrightarrow I^{p,\bullet}.
\]
のコホモロジーをとることで得られる。

\VerdierSemanticAnchor{ch4d-4-6-8-proof-end}
ところがこの射は擬同型である。これで命題が示された。

\VerdierSemanticAnchor{ch4d-4-6-9}
\medskip\noindent\textbf{4.6.9.}\quad
$F\colon D(\mathcal A)\longrightarrow\mathcal B$ をコホモロジー的関手、
$Y$ を $\operatorname{comp}(\mathcal A)$ の対象とし、
$Y\longrightarrow I^{\bullet\bullet}$ を Cartan--Eilenberg 入射分解とする。
関手 $F$ は、対応する Cartan--Eilenberg スペクトル対象 (4.6.7) を
$\mathcal B$ に値をとるスペクトル対象へ移し、したがって
\emph{Cartan--Eilenberg スペクトル系列}と呼ばれるスペクトル系列を与える。
このスペクトル系列の項を $(\mathrm{C.E.})_r^{p,q}$ と記す
（第 II 章、4.3.3.1）。その初期項は (4.1.9.1)

\VerdierSemanticAnchor{ch4d-4-6-9-initial-terms}
\[
(\mathrm{C.E.})_1^{p,q}
=F^{p+q}\bigl(I^{\bullet,p}[-p]\bigr)
=F^q(I^{\bullet,p}),
\]
\endgroup
